\section{Three Finite-Message Effects}
\label{sec:audit}

\subsection{Finite Input Is Not an Infinite Random Stream}

An arithmetic decoder can simulate a token source when its driving bits have the
required random distribution. Finishing a finite message changes that premise.
The original NLS implementation~\cite{ziegler2019neural} maintains a fixed-precision
interval and renormalizes shared leading bits. When its lookahead extends beyond
the message, it fills the missing positions with zeros. Its reported KL compares
rounded token frequencies with model probabilities; this calculation does not
account for the deterministic suffix of the actual input.

Our harness extracts and executes the original encoder and its utility functions
unchanged from a pinned 2019 revision. A synthetic model assigns probability
$1/64$ to each of 64 active tokens. We enumerate all 256 eight-bit messages,
with precision 16, candidate cap 64, and sentence completion disabled. Each
token has an injective synthetic Unicode rendering, and rendering followed by
tokenization is checked for every output. Each output contains two tokens.
The first carries six bits; the remaining two input
bits and four zero padding bits restrict the final token to four choices.
Consequently,
\begin{equation}
 \operatorname{TV}(R_{\mathrm{last}},U_{64})=\frac{15}{16},
 \qquad \widehat D_{\mathrm{KL}}=0.
\end{equation}
The same counting argument applies to $B=6k+2$, including 128-bit messages.
This is an algorithm-level counterexample to the diagnostic, not an empirical
claim about GPT-2 output quality. It holds with a uniform finite input, so
normalizing plaintext through encryption does not remove this padding effect.
The injective text channel has 256 equally likely stego strings among 4096
possible natural two-symbol strings, giving complete-text TV $15/16$ as well.
Meteor masks its lookahead with fresh pseudorandom bits after padding; the
specific argument above does not apply to that construction~\cite{kaptchuk2021language}.

\subsection{Integer-Interval Sampling Law}

Ghostext's legacy codec uses a different representation. It interprets an entire
$n$-byte segment as an integer in $[0,2^{8n})$ and narrows an integer interval
until only one integer remains. Suppose the current width is $W$, total
frequency is $F$, and token $i$ has cumulative integer frequencies $c_i,c_{i+1}$.
Even under a uniform target within the current interval, its selection
probability is
\begin{equation}
 R_W(i)=\frac{\lfloor Wc_{i+1}/F\rfloor-\lfloor Wc_i/F\rfloor}{W}.
 \label{eq:integer-law}
\end{equation}
This law need not equal $Q(i)=(c_{i+1}-c_i)/F$. Each endpoint floor contributes
less than one unit of count error, giving
\begin{equation}
 \operatorname{TV}(R_W,Q)<\frac{N}{2W},
 \label{eq:local-rounding}
\end{equation}
for $N$ candidate tokens. The bound can be large when $W$ is small.

An exhaustive uniform-64 example again gives TV $15/16$, but for a different
reason. With an eight-bit target and $F=4096$, the first token leaves width four.
Only four of the next 64 subintervals are nonempty. Of the 4096 possible natural
two-token sequences, only 256 are admissible. This effect exists with exact
frequency quantization and an identity tokenizer.

It also gives a keyless test for a framed protocol. Ghostext's bootstrap length
is public, so an observer can initialize its interval, replay candidate tables,
and reject any sequence that enters a zero-width subinterval. No authentication
key is needed. We retain separate reasons for rejection: candidate exclusion,
violation of a deterministic no-payload step, and zero-width interval selection.
A positive result means only that the prefix is admissible, not that it contains
an authenticated packet.

\subsection{Other Distribution Changes}

The released quantizer reserves one frequency for every candidate before
allocating the remaining $F-N$ counts. An $N/(2F)$ error claim therefore cannot
be borrowed from ordinary rounding. With $N=64$, $F=4096$, and probabilities
approximately $(1,0,\ldots,0)$, its measured TV is $0.015381$, exceeding
$N/(2F)=0.0078125$. Our revised audit computes
$\frac12\sum_i|p_i-f_i/F|$ directly and measures interval rounding separately
using Eq.~\eqref{eq:integer-law}.

Truncation, removal of retokenization-unstable candidates, greedy emission at
low entropy, successful-attempt selection, and the final stopping rule can all
change the visible distribution. In particular, fresh retries do not remove
selection bias. For an attempt law $\mu$ and success event $A$, the first
successful independent attempt has law $\mu(\cdot\mid A)$, regardless of the
retry cap conditional on eventual success. At the attempt-trajectory level,
$\operatorname{TV}(\mu,\mu(\cdot\mid A))=1-\mu(A)$; projecting to visible text
can reduce this distance, but cannot justify ignoring the conditioning.
